\chapter{Safety}
\label{ch:safety}

\begin{wbwarning}
Read this chapter before operating the machine. The bond head moves under power
and without warning, and the machine can begin moving as soon as it is switched
on.
\end{wbwarning}

\section{The machine moves at power-up}

Switching the machine on starts an automatic homing sequence. The Y stage
travels to find its limit switch, then the bond head initialises. This happens
without any operator action.

\begin{wbwarning}
Keep hands, tools and workpieces clear of the stage and bond head while the
machine is powering up. Do not switch it on with anything resting on the stage.
\end{wbwarning}

\section{Pinch and crush points}

The bond head descends under force onto the workholder, and the Y stage travels
beneath it. Fingers placed between them will be trapped.

\begin{wbwarning}
Never place any part of your hand between the tool and the workholder while the
machine is powered. The bond head can descend at up to 10\,mm/s under a load of
up to 150\,grams, triggered by a single mouse-button press.
\end{wbwarning}

\section{\keycap{RESET} stops the machine, but not gently}

\keycap{RESET} restarts the machine immediately, from any state, including
mid-weld. Motion stops wherever it is and the machine re-homes.

\begin{wbwarning}
\keycap{RESET} is not an emergency stop and must not be relied on as one. It
does not remove power from the motors --- it restarts the controller, which
then moves the axes again to re-home. If you need to make the machine safe,
remove its power.
\end{wbwarning}

\TODO{State what the actual emergency stop provision is --- a mains switch, an
E-stop mushroom, an isolator --- and where it is located. Add a labelled photo.
If there is no E-stop, say so explicitly and state that the mains disconnect is
the means of isolation.}

\section{Ultrasonic energy}

The tool vibrates at approximately 60\,kHz during a weld. This is above the
range of human hearing but the energy is real.

\begin{wbwarning}
Do not touch the tool, the transducer, or the horn while a weld is in progress.
Ultrasonic contact with skin causes burns and tissue damage without an
immediate pain warning.
\end{wbwarning}

The transducer and horn also become warm in continuous use. Allow them to cool
before handling.

\section{Lighting}

The machine has an area light and a spotlight.

\TODO{Confirm whether the spotlight is a laser. If it is, state its class and
add the required warning against looking into the beam or viewing it with
optical instruments, along with the label location. If it is an LED, say so and
remove this block.}

\section{Electrical}

\TODO{Supply the electrical safety content: supply voltage and current, earthing
requirements, which panels may be opened by an operator (normally none), the
isolation procedure before any cover is removed, and fuse ratings and
locations.}

\section{Workpiece and consumables}

\TODO{Supply handling guidance for the materials in use --- wire alloys, any
solvents used for tool cleaning, and ESD precautions for the devices being
bonded.}

\section{Regulatory}

\TODO{Supply declarations of conformity, applicable standards, CE/UKCA or other
markings, and the location of the machine's rating plate.}
